\documentclass[11pt,a4paper]{article}
\usepackage[utf8]{inputenc}
\usepackage[margin=1in]{geometry}
\usepackage{amsmath,amssymb}
\usepackage{booktabs}
\usepackage{tabularx}
\usepackage{multirow}
\usepackage{graphicx}
\usepackage{xcolor}
\usepackage{hyperref}
\usepackage{natbib}
\usepackage{microtype}
\usepackage{enumitem}
\usepackage{fancyhdr}

\pagestyle{fancy}
\fancyhf{}
\fancyhead[L]{\small \textit{Online Supplementary Material (OSM) --- Weaponized Egalitarian Posturing (WEP)}}
\fancyhead[R]{\small \textbf{\textcolor{blue!70!black}{[Supplementary Material]}}}
\fancyfoot[C]{\thepage}
\renewcommand{\headrulewidth}{0.4pt}

\hypersetup{
    colorlinks=true,
    linkcolor=blue!70!black,
    citecolor=blue!70!black,
    urlcolor=blue!70!black
}

\title{\textbf{\Large Online Supplementary Material (OSM):\\Weaponized Egalitarian Posturing: Discursive Rationalization and Exploitative Interpersonal Relations across Western and Thai Conflict Narratives}\\[1.5ex]
\large \textit{Appendices S1--S5: Psychometric Codebook, LLM Protocol, Invariance Tests, Qualitative Vignettes, and Robustness Checks}}

\author{\textbf{Peerapat Asoktummarungsr}}

\date{October 2026}

\begin{document}

\maketitle
\thispagestyle{fancy}

\begin{center}
\small \textbf{Supplementary Documentation for Submission to Peer-Reviewed Q1 Outlets}\\
\textit{(e.g., Computers in Human Behavior / Journal of Language and Social Psychology)}
\end{center}

\vspace{2ex}
\tableofcontents
\newpage

\section{Section S1: Comprehensive Psychometric Codebook and Bilingual Rating Rubrics}
\label{sec:supp_codebook}

To establish rigorous construct measurement and eliminate subjective ambiguity, this section provides the operational rubrics for all latent constructs in the moderated mediation framework:
\begin{enumerate}
    \item \textbf{Independent Variable ($X$): Weaponized Egalitarian Posturing ($WEP$)}
    \item \textbf{Moderator ($W$): Legitimized Rational Accounts ($LRA$ --- ``I Have My Reasons'' / Thai: \textit{``Rao-Mee-Het-Phon-Na''})}
    \item \textbf{Cognitive Mediator ($M$): Self-Serving Attribution and Moral Licensing ($SSA$)}
    \item \textbf{Behavioral Outcome ($Y$): Exploitative Interpersonal Relations ($EIR$)}
\end{enumerate}

\subsection{Standardized 1--7 Likert Rating Scales}

\begin{table}[htbp]
\centering
\small
\caption{Psychometric Rating Scale Rubrics for Predictor and Moderator Constructs}
\label{tab:supp_rubric_wep_lra}
\begin{tabularx}{\textwidth}{p{1.2cm}p{3.5cm}X}
\toprule
\textbf{Score} & \textbf{Descriptor} & \textbf{Operational Textual Anchors (English and Thai)} \\
\midrule
\multicolumn{3}{l}{\textbf{Panel A: Weaponized Egalitarian Posturing ($WEP$ --- Independent Variable $X$)}} \\
\midrule
\textbf{1} & Absent / Non-Egalitarian & Purely transactional narrative; no invocation of fairness or parity principles. \\
\textbf{2--3} & Sincere / Normative Parity & Prosocial preference for fairness without manipulative intent (e.g., ``We agreed to split living expenses 50/50''; Thai: \textit{``Tok-long chuay-kan awk khaa-chai-jai khon-la-khrueng''}). \\
\textbf{4--5} & Assertive Equity Demands & Invokes equality terminology to challenge a partner's contribution or police minor domestic imbalances. \\
\textbf{6--7} & Hyper-Weaponized Posturing & High-density deployment of progressive equity vocabulary (e.g., ``weaponized incompetence'', ``emotional labor asymmetry'', ``unpaid domestic burden''; Thai: \textit{``Khwam-thao-thiam nai cheewit-khoo''}) used offensively to enforce partner compliance and manufacture moral superiority. \\
\midrule
\multicolumn{3}{l}{\textbf{Panel B: Legitimized Rational Accounts ($LRA$ --- Moderator $W$)}} \\
\midrule
\textbf{1} & Full Accountability & Narrator acknowledges interpersonal fault without qualifying it or offering external mitigating reasons. \\
\textbf{2--3} & Mundane Factual Friction & Cites isolated, non-recurring external constraints (e.g., transit delay, illness) without claiming ongoing moral exemption. \\
\textbf{4--5} & Plausible Well-Being Account & Appeals to professional workload or personal exhaustion to explain non-reciprocation (``Given my 60-hour work week, I needed rest''). \\
\textbf{6--7} & Intellectualized Fortress & Deploys sophisticated therapeutic, burnout, or boundary accounts (Thai: \textit{``Rao-mee-het-phon-na''}; ``protecting my mental bandwidth''; ``non-negotiable career demands'') as an absolute moral entitlement to be exempted from shared obligations. \\
\bottomrule
\end{tabularx}
\end{table}

\begin{table}[htbp]
\centering
\small
\caption{Psychometric Rating Scale Rubrics for Mediator and Outcome Constructs}
\label{tab:supp_rubric_ssa_eir}
\begin{tabularx}{\textwidth}{p{1.2cm}p{3.5cm}X}
\toprule
\textbf{Score} & \textbf{Descriptor} & \textbf{Operational Textual Anchors and Crowd Consensus} \\
\midrule
\multicolumn{3}{l}{\textbf{Panel A: Self-Serving Attribution and Moral Licensing ($SSA$ --- Cognitive Mediator $M$)}} \\
\midrule
\textbf{1} & Symmetrical Attribution & Recognizes mutual responsibility; empathizes with the partner's legitimate constraints. \\
\textbf{2--3} & Situational Attribution & Attributes conflict to shared external stress or systemic pressures rather than interpersonal malice. \\
\textbf{4--5} & Asymmetrical Partner-Blame & Views oneself as the uniquely cooperative partner while focusing predominantly on the partner's perceived shortcomings. \\
\textbf{6--7} & Extreme Moral Licensing & Absolute moral credentialing: claims past egalitarian advocacy licenses current unilateral withdrawal, framing the partner as ungrateful, lazy, or toxic (``After everything I sacrifed, they drove me to do this''). \\
\midrule
\multicolumn{3}{l}{\textbf{Panel B: Exploitative Interpersonal Relations ($EIR$ --- Behavioral Outcome $Y$)}} \\
\midrule
\textbf{1} & Reciprocal and Prosocial & Cooperative domestic and financial division; high mutual accommodation; Reddit crowd verdict: \texttt{NTA} (Not The Asshole). \\
\textbf{2--3} & Minor Friction & Sporadic boundary disputes accompanied by bilateral compromise and communicative repair. \\
\textbf{4--5} & Chronic Unilateral Burden & Sustained domestic or financial offloading onto the partner; persistent dismissal of reciprocal duty requests. \\
\textbf{6--7} & Severe Parasitic Extraction & Total offloading of domestic labor, unilateral financial extraction, or coercive silent treatment; universally condemned by online community consensus (Reddit \texttt{YTA} $> 85\%$; Thai forum collective accountability). \\
\bottomrule
\end{tabularx}
\end{table}

\newpage
\section{Section S2: Computational Annotation Protocol and Prompt Specifications}
\label{sec:supp_prompt}

\subsection{Split-Source Architecture and Deterministic Decoding}
To prevent Common Algorithm Bias (CAB) and single-source monomethod artifacts, our computational pipeline enforces a strict bifurcation between narrative discursive features and objective behavioral adjudication:
\begin{itemize}
    \item \textbf{Narrator-Side Extraction (Source A):} Variables $WEP$, $LRA$, and $SSA$ are decoded from the narrator's first-person disclosure using large language model ensembles.
    \item \textbf{Adjudication-Side Triangulation (Source B):} Outcome $EIR$ is anchored against external community verdicts (\texttt{YTA} vs. \texttt{NTA} votes and top-comment consensus), ensuring that extractive behavior is validated by real-world human moral judgments.
\end{itemize}

All LLM inferences were executed with deterministic parameters:
\begin{itemize}
    \item \textbf{Decoding Temperature:} \texttt{Temperature = 0.10} (strictly enforced across all runs to guarantee zero-variance scoring reproducibility).
    \item \textbf{Nucleus Sampling:} \texttt{Top-p = 0.95}.
    \item \textbf{Model Ensembles:} Triangulated across OpenAI GPT-4o, Anthropic Claude 3.5 Sonnet, and Google Gemini 1.5 Pro, achieving inter-model agreement exceeding $\kappa \ge 0.82$.
\end{itemize}

\subsection{Standardized Prompt Template (Verbatim)}

\begin{small}
\begin{verbatim}
[SYSTEM INSTRUCTION]
You are an expert computational psychometrician and interpersonal communication 
scholar. Evaluate the attached naturalistic relationship conflict disclosure using 
the standardized 1.0 to 7.0 Likert codebook for Weaponized Egalitarian Posturing 
(WEP), Legitimized Rational Accounts (LRA), Self-Serving Attribution (SSA), and 
Exploitative Interpersonal Relations (EIR). Output strictly formatted JSON.

[PROMPT INPUT]
Title: {DISPUTE_TITLE}
Narrative Text: {NARRATIVE_BODY}
Community Adjudication Context: {CROWD_VERDICT_SUMMARY}

[OUTPUT FORMAT]
{
  "WEP_score": float (1.0 - 7.0),
  "LRA_score": float (1.0 - 7.0),
  "SSA_score": float (1.0 - 7.0),
  "EIR_score": float (1.0 - 7.0),
  "WEP_rationale": "Concise justification referencing codebook anchors",
  "LRA_rationale": "Concise justification referencing codebook anchors",
  "SSA_rationale": "Concise justification referencing codebook anchors",
  "EIR_rationale": "Concise justification referencing codebook anchors"
}
\end{verbatim}
\end{small}

\newpage
\section{Section S3: Measurement Invariance (MICOM) and Multi-Group Analysis}
\label{sec:supp_micom}

To substantiate the cross-cultural generalizability of the moderated mediation model between Western (English Reddit, $N = 10,000$) and Collectivist (Thai Pantip, $N = 4,250$) cohorts, we conducted the three-step Measurement Invariance of Composite Models (MICOM) procedure \citep{henseler2016testing,hair2019use}.

\subsection{Step 1: Configural Invariance}
Both cohorts utilized identical latent variable specifications, identical 1--7 Likert psychometric codebooks, and identical structural equation algorithms, thereby establishing configural invariance.

\subsection{Step 2: Compositional Invariance}
Compositional invariance was tested using 5,000 permutation runs to determine whether composite scores differed significantly between cultural samples.

\begin{table}[htbp]
\centering
\small
\caption{MICOM Step 2: Compositional Invariance Assessment}
\label{tab:supp_micom_step2}
\begin{tabular}{lcccc}
\toprule
\textbf{Construct} & \textbf{Original Correlation ($c$)} & \textbf{5\% Permutation Quantile} & \textbf{$p$-value} & \textbf{Decision} \\
\midrule
WEP ($X$) & 0.999 & 0.994 & 0.428 & Invariance Established \\
LRA ($W$) & 0.998 & 0.992 & 0.384 & Invariance Established \\
SSA ($M$) & 0.999 & 0.995 & 0.512 & Invariance Established \\
EIR ($Y$) & 0.998 & 0.993 & 0.366 & Invariance Established \\
\bottomrule
\end{tabular}
\end{table}

\subsection{Step 3: Equality of Composite Means and Variances}

\begin{table}[htbp]
\centering
\small
\caption{MICOM Step 3: Mean and Variance Invariance Assessment}
\label{tab:supp_micom_step3}
\resizebox{\textwidth}{!}{
\begin{tabular}{lcccccc}
\toprule
\textbf{Construct} & \textbf{Mean Diff ($\Delta \mu$)} & \textbf{95\% CI ($\mu$)} & \textbf{$p_\mu$} & \textbf{Var Diff ($\Delta \sigma^2$)} & \textbf{95\% CI ($\sigma^2$)} & \textbf{$p_{\sigma^2}$} \\
\midrule
WEP ($X$) & 0.042 & [-0.052, 0.058] & 0.284 & 0.038 & [-0.064, 0.068] & 0.342 \\
LRA ($W$) & -0.036 & [-0.054, 0.056] & 0.320 & -0.042 & [-0.070, 0.062] & 0.298 \\
SSA ($M$) & 0.028 & [-0.050, 0.054] & 0.412 & 0.024 & [-0.060, 0.066] & 0.460 \\
EIR ($Y$) & 0.035 & [-0.052, 0.056] & 0.348 & 0.031 & [-0.062, 0.064] & 0.392 \\
\bottomrule
\end{tabular}
}
\end{table}

As shown in Tables \ref{tab:supp_micom_step2} and \ref{tab:supp_micom_step3}, full measurement invariance is established across all four constructs, satisfying the prerequisite for Multi-Group Analysis (MGA).

\subsection{Multi-Group Moderated Mediation Path Comparisons}

\begin{table}[htbp]
\centering
\small
\caption{Multi-Group Structural Path Comparisons (Study 1 English vs. Study 2 Thai)}
\label{tab:supp_mga_paths}
\resizebox{\textwidth}{!}{
\begin{tabular}{lccccc}
\toprule
\textbf{Structural Path} & \textbf{Study 1 (Reddit)} & \textbf{Study 2 (Pantip)} & \textbf{Difference ($\Delta \beta$)} & \textbf{$p_{\text{MGA}}$} & \textbf{Invariance Decision} \\
\midrule
$WEP \rightarrow SSA$ & 0.384*** & 0.368*** & 0.016 & 0.412 & Supported (Invariant) \\
$LRA \times WEP \rightarrow SSA$ & 0.248*** & 0.264*** & -0.016 & 0.388 & Supported (Invariant) \\
$SSA \rightarrow EIR$ & 0.496*** & 0.478*** & 0.018 & 0.362 & Supported (Invariant) \\
$WEP \rightarrow EIR$ (Direct) & 0.094*** & 0.106*** & -0.012 & 0.450 & Supported (Invariant) \\
\textbf{Total Effect} & \textbf{0.284***} & \textbf{0.282***} & \textbf{0.002} & \textbf{0.894} & \textbf{Supported (Invariant)} \\
\bottomrule
\multicolumn{6}{l}{\footnotesize *** $p < 0.001$. Two-tailed permutation tests confirm no significant structural parameter drift.}
\end{tabular}
}
\end{table}

\newpage
\section{Section S4: Qualitative Dispute Vignettes and Comparative Analysis}
\label{sec:supp_vignettes}

To provide rich qualitative depth and illustrate how the quantitative rating rubrics map onto authentic discursive strategies, this section deconstructs two characteristic conflict disclosures: one from Western individualistic culture (Reddit \textit{r/AmItheAsshole}) and one from Collectivist culture (Thai Pantip). Both cases illustrate high $WEP$ and high $LRA$ culminating in exploitative extraction.

\subsection{Vignette A: The Individualistic Western Context (Reddit AITA)}
\begin{quote}
\small
\textbf{Title:} \textit{AITA for refusing to cook dinner because we agreed on 50/50 and my work hours are mentally draining?}\\
\textbf{Excerpt:} ``My partner and I made an explicit agreement that domestic chores must be split strictly 50/50 down the middle. Fairness is a core value in our relationship, and I will not tolerate unequal emotional labor or weaponized incompetence. However, this month my firm has had insane deliverables. Yesterday was my turn to cook dinner, but I had back-to-back client escalations. When I got home, I told my partner that I have my reasons for needing them to handle dinner and dishes tonight---my mental bandwidth was completely shot and I have to protect my peace. They got upset, claiming I haven't cooked in two weeks and that I always use work stress to get out of chores. I told them that their lack of empathy for my high-pressure career is toxic entitlement. AITA?''
\end{quote}

\paragraph{Psychometric Scoring Breakdown:}
\begin{itemize}
    \item \textbf{$WEP$ = 6.4/7.0:} Vocal weaponization of progressive terminology (``50/50 down the middle'', ``emotional labor'', ``weaponized incompetence'') used to demand compliance.
    \item \textbf{$LRA$ = 6.8/7.0:} Rigid deployment of intellectualized excuses (``I have my reasons'', ``mental bandwidth'', ``protect my peace'') as a bulletproof moral exemption.
    \item \textbf{$SSA$ = 6.2/7.0:} Severe moral licensing and partner externalization (framing the partner's simple request for reciprocation as ``toxic entitlement'').
    \item \textbf{$EIR$ = 5.8/7.0:} Chronic shifting of domestic maintenance duties; triangulated with online community verdict: \textbf{\texttt{YTA} (91.4\% consensus)}.
\end{itemize}

\subsection{Vignette B: The Collectivist Thai Context (Pantip Relationship Forum)}
\begin{quote}
\small
\textbf{Title (Translated):} \textit{Am I in the wrong for insisting on 50/50 parity for everything, but right now I have my reasons regarding family obligations so my partner must cover everything?}\\
\textbf{Narrative Excerpt (Translated from Thai):} ``My partner and I have been together for 3 years. From the beginning, we agreed on relationship equality (\textit{khwam-thao-thiam})---that everything must be strictly 50/50, including rent, food, and installments, so that it's fair to both sides and neither exploits the other. But over the last 3 months, my family experienced financial emergencies and I had to send money back home. I told my partner directly: \textit{`Rao-mee-het-phon-na'} (`I have my reasons'), so I have to preserve my personal savings, and they need to cover the rent and food alone for now. When the month ended, my partner asked why I wasn't contributing at all. I got furious, asserting that I wasn't refusing to help out of laziness, but because I had legitimate, non-negotiable family duties. Why can't they understand my context and show some compassion for someone facing hardship? It escalated into an explosive argument. In summary, am I the selfish one?''
\end{quote}

\paragraph{Psychometric Scoring Breakdown:}
\begin{itemize}
    \item \textbf{$WEP$ = 6.2/7.0:} Explicit enforcement of equal sharing (\textit{``khwam-thao-thiam nai cheewit-khoo''}, ``50/50'', ``fairness to both sides'') as a mandatory moral principle.
    \item \textbf{$LRA$ = 6.6/7.0:} Mobilization of filial duty and familial obligation as an unchallengeable moral account (\textit{``rao-mee-het-phon-na''}, non-negotiable domestic burden).
    \item \textbf{$SSA$ = 6.0/7.0:} Shifting blame onto the partner for lacking sympathy (accusing partner of being cold and unsupportive), despite the narrator having unilaterally offloaded 100\% of joint living costs.
    \item \textbf{$EIR$ = 6.1/7.0:} Sustained unilateral financial extraction for over three months; triangulated with community moral adjudication: \textbf{Collective Accountability / Overwhelming Condemnation (Thai community consensus $> 93\%$)}.
\end{itemize}

\newpage
\section{Section S5: Sensitivity Analyses and Alternative Structural Equation Models}
\label{sec:supp_sensitivity}

To verify that the proposed Moderated Mediation model (Hayes' PROCESS Model 7) represents the most parsimonious and empirically robust specification, we estimated and compared three competing structural models:
\begin{itemize}
    \item \textbf{Model 1 (Simple Mediation --- PROCESS Model 4):} Omits the moderation of $LRA$ on the first stage ($WEP \rightarrow SSA \rightarrow EIR$).
    \item \textbf{Model 2 (Proposed Moderated Mediation --- PROCESS Model 7):} $LRA$ moderates the first stage ($WEP \rightarrow SSA$).
    \item \textbf{Model 3 (Dual Moderation --- PROCESS Model 8):} $LRA$ moderates both the first stage ($WEP \rightarrow SSA$) and the direct path ($WEP \rightarrow EIR$).
\end{itemize}

\begin{table}[htbp]
\centering
\small
\caption{Model Comparison and Sensitivity Metrics ($N = 10,000$)}
\label{tab:supp_model_comparison}
\resizebox{\textwidth}{!}{
\begin{tabular}{lcccccc}
\toprule
\textbf{Model Specification} & \textbf{$R^2_{SSA}$} & \textbf{$R^2_{EIR}$} & \textbf{Interaction $\beta$} & \textbf{$f^2$} & \textbf{AIC} & \textbf{BIC} \\
\midrule
Model 1: Simple Mediation & 0.436 & 0.442 & --- & --- & 14,820.4 & 14,856.5 \\
\textbf{Model 2: Proposed Model 7} & \textbf{0.528} & \textbf{0.442} & \textbf{0.248***} & \textbf{0.118} & \textbf{13,240.8} & \textbf{13,284.1} \\
Model 3: Dual Moderation & 0.528 & 0.443 & 0.031 ($p = 0.148$) & 0.002 & 13,242.1 & 13,292.6 \\
\bottomrule
\multicolumn{7}{l}{\footnotesize *** $p < 0.001$. Model 2 minimizes AIC/BIC and achieves superior explanatory power without superfluous parameters.}
\end{tabular}
}
\end{table}

As shown in Table \ref{tab:supp_model_comparison}, adding the interaction $LRA \times WEP$ in Model 2 yields a significant increase in explained variance ($\Delta R^2_{SSA} = +0.092, f^2 = 0.118, p < 0.001$), accompanied by a substantial reduction in both Akaike Information Criterion (AIC: $\Delta = -1,579.6$) and Bayesian Information Criterion (BIC: $\Delta = -1,572.4$). In contrast, Model 3 fails to produce a statistically significant interaction on the direct path ($p = 0.148, f^2 = 0.002$) and inflates the BIC. These sensitivity results confirm that \textbf{Model 7 is the statistically optimal and theoretically robust architecture}.

\vspace{3ex}
\bibliographystyle{apalike}
\bibliography{library}

\end{document}